\section{Remaining classification and induction problems}\label{sec:open}

Three questions remain beyond the proved range.
\begin{enumerate}
\item The necessity direction of \cref{conj:odd} is open from rank eleven onward. A uniform argument must force the adjacent-pair relations and the ordered Boolean chain without a separate elimination in each rank.
\item The companion indices $7\le k\le m-1$ are unproved, beginning with $k=7$ in rank seventeen. The boundary inverse is known at every index (\cref{lem:inverse,thm:folding}); the missing step is the reduction on the Nahm side.
\item In even rank, sufficiency is open in rank six for five of the nine pairs of \cref{thm:r6}, and both directions are open from rank eight onward. \Cref{conj:even} below gives the expected answer, and \cref{conj:evenwron} reduces sufficiency to explicit theta-determinant identities, as \cref{conj:wron} does in odd rank. A related question is whether a number of cusp conditions independent of the rank suffices; ranks two, three, five, six, seven and nine need three coefficients, while rank four needs four.
\end{enumerate}

\subsection{Even ranks}
Ranks two, four and six suggest the following description.

\begin{conjecture}[Even-rank classification]\label{conj:even}
Let $r=2m$ with $m\ge1$. Then $\M_r$ contains the following $2m+3$ pairs, and for $m\ne2$ it consists exactly of them.
\begin{enumerate}
\item For integers $0\le x_0\le x_1\le\cdots\le x_{m-1}\le1$ with $k=\sum_ix_i$,
\[
B=(x_0,\,-x_1,\,x_1,\,\dots,\,-x_{m-1},\,x_{m-1},\,0),\qquad C=-\frac{r(2r+1)}{48(r+2)}+\frac{k(r+1-k)}{2(r+2)}.
\]
\item For integers $1\ge y_1\ge\cdots\ge y_{m-1}\ge0$ with $j=\sum_iy_i$,
\[
B=(y_1,\,-y_1,\,\dots,\,y_{m-1},\,-y_{m-1},\,0,\,\tfrac12),\qquad C=-\frac{r(r-4)}{24(r+2)}+\frac{j(r-j)}{2(r+2)}.
\]
\item $B=(1,-1,\dots,1,-1,1,-\tfrac12)$, with $m-1$ pairs $(1,-1)$, and $C=r/12$.
\item The vector of (iii) with $(1,-1)$ added to its first two coordinates, and $C=C_{m-1}+1$, where $C_{m-1}$ is the scalar in (ii) with $j=m-1$.
\end{enumerate}
For $m=2$, \cref{thm:r4} shows that $\M_4$ consists of these seven pairs together with the exceptional pair $\bigl((\frac32,-\frac32,\frac12,0),\frac9{16}\bigr)$.
\end{conjecture}

The sector (i) has the same scalar as \cref{conj:odd} with $z=0$. \Cref{thm:r2,thm:r4,thm:r6} agree with the conjecture in ranks two, four and six; in rank six only necessity and two of the nine sufficiency statements are proved. In ranks two and six the pair (iv) satisfies $F_{r,B_{\mathrm{(iv)}},0}=2q^{-1}F_{r,B_{\mathrm{(ii)},j=m-1},0}$, exactly in rank two and to order $q^{24}$ in rank six, which explains the shift of $C$ by one.

\subsection{An even-rank Wronskian formula}
The sums in \cref{conj:even} appear to satisfy theta-determinant identities parallel to \eqref{eq:wron}, with an alternating character on the lattice and without the factor $\prod_iv_i$. Let $r=2m$, $L=2m+2$ and $\nu=m(2m-1)$, and put
\[
\chi_0(v)=\prod_{1\le i<j\le m}(v_i^2-v_j^2),\qquad \epsilon_\tau(v)=(-1)^{\sum_i(v_i-\tau_i)/L}\quad(v\in\tau+L\Z^m).
\]
Define increasing tuples
\begin{align*}
\tau^{(0)}_k&=(\tfrac12,\tfrac32,\dots,m+\tfrac12)\setminus\{m+\tfrac12-k\}, & \tau^{(0)}_*&=(\tfrac12,\dots,m-\tfrac12),\\
\tau^{(1/2)}_j&=(0,1,\dots,m)\setminus\{m-j\}, & \tau^{(1/2)}_*&=(0,1,\dots,m-1),
\end{align*}
for $0\le k,j\le m$, and let
\[
R^{(z)}_\tau(q)=\frac{(-q^{1/2+z};q)_\infty^r}{(q;q)_\infty^\nu}\sum_{v\in\tau+L\Z^m}\epsilon_\tau(v)\,\frac{\chi_0(v)}{\chi_0(\tau^{(z)}_*)}\,q^{\frac{\|v\|^2-\|\tau\|^2}{2L}}.
\]

\begin{conjecture}[Even-rank Wronskian identity]\label{conj:evenwron}
Let $B^{\mathrm{(i)}}_k$ be the vector of \cref{conj:even}(i) with $\sum_ix_i=k$, let $B^{\mathrm{(ii)}}_j$ ($0\le j<m$) be the vector of \cref{conj:even}(ii) with $\sum_iy_i=j$, and let $B^{\mathrm{(ii)}}_m$ be the vector of \cref{conj:even}(iii). Then
\[
F_{2m,B^{\mathrm{(i)}}_k,0}=R^{(0)}_{\tau^{(0)}_k},\qquad F_{2m,B^{\mathrm{(ii)}}_j,0}=c_j\,R^{(1/2)}_{\tau^{(1/2)}_j},\qquad c_j=\begin{cases}1,&j<m,\\2,&j=m.\end{cases}
\]
Moreover, the vector $B^{\mathrm{(iv)}}$ of \cref{conj:even}(iv) satisfies $F_{2m,B^{\mathrm{(iv)}},0}=2q^{-1}F_{2m,B^{\mathrm{(ii)}}_{m-1},0}$.
\end{conjecture}

The cases $B^{\mathrm{(i)}}_0=0$ and $B^{\mathrm{(ii)}}_0=(0,\dots,0,\frac12)$ are the principal and last-coordinate half-shift formulas of \cite{CTW26} for even rank; the alternating character is their sign $\varepsilon=(-1)^{r+1}$. In rank two the last relation is proved in \cref{thm:r2}. In odd rank the analogous relation \eqref{eq:halfshift} follows from the simultaneous shift of \cref{lem:simshift}, which comes from the odd-rank reduction of Shi and Wang. Their even-rank reduction \cite[(1.10)]{SW26}, which holds with $x_0=x_1^{-1}$, admits no shift of this kind: the factors $x_0^{-n_1}$ and $(x_2/x_0)^{n_2}$ leave an uncancelled term $q^{hn_2}$. A proof of the last relation, and of \cref{conj:evenwron} beyond the principal cases, therefore seems to require the Hall--Littlewood identities of \cite{BW15} at non-principal specializations, as in \cref{sec:theta}.

\begin{proposition}\label{prop:evenwron}
For each tuple $\tau$ above, $q^{C}R^{(z)}_\tau$ is a modular function on a congruence subgroup, where $C$ is the scalar attached to the corresponding pair in \cref{conj:even}. Consequently, \cref{conj:evenwron} implies that every pair of \cref{conj:even} is modular, and in particular that $\M_6$ consists of exactly the nine pairs of \cref{thm:r6}.
\end{proposition}

\begin{proof}
For rational $a$, the series $\theta^-_{a,L}(q)=\sum_{v\in a+L\Z}(-1)^{(v-a)/L}q^{v^2/(2L)}$ has weight $1/2$ on a congruence subgroup. As in the proof of \cref{prop:rhs-modular}, expanding the Wronskian $\det(\mathcal D^{i-1}\theta^-_{\tau_j,L})_{i,j=1}^m$ and using the Vandermonde determinant in the variables $v_j^2$ gives a constant multiple of $\sum_{v\in\tau+L\Z^m}\epsilon_\tau(v)\chi_0(v)q^{\|v\|^2/(2L)}$, which has weight $m/2+m(m-1)=\nu/2$. The eta-quotient factors are those of \cref{prop:rhs-modular}, and the weights cancel. The power of $q$ that normalizes the product is
\[
\frac{\|\tau\|^2}{2L}-\frac\nu{24}-\begin{cases}\frac r{48},&z=0,\\[2pt]-\frac r{24},&z=\frac12,\end{cases}
\]
and direct evaluation of $\|\tau\|^2$ shows that it equals the scalar of \cref{conj:even} in sectors (i)--(iii). The relation for $B^{\mathrm{(iv)}}$ multiplies by $2q^{-1}$, which shifts the scalar by one, as in (iv). Together with \cref{thm:r6}, this gives the last assertion.
\end{proof}

\subsection{Even rank: the first companions}\label{subsec:evenproof}
The argument of \cref{sec:theta,sec:boundary,sec:cert} transfers to the sector (i) of \cref{conj:even}, with the type-C identity of Bartlett and Warnaar replaced by its type-B companion. Throughout, $r=2m$, $L=2m+2$, $\nu=m(2m-1)$ and $w=q^{1/2}$. For generic $\hat x=(x_1,\dots,x_{m-1})$ put $x=(\hat x,1)$; thus $x_m=1$ is frozen. Define
\[
\HL^B_m(t;\hat x;q)=\sum_{\lambda:\ \lambda_1\le2}q^{|\lambda|/2}\,t^{m_1(\lambda)}\,P'_\lambda(x_1,x_1^{-1},\dots,x_{m-1},x_{m-1}^{-1},1;q),
\]
a sum over Hall--Littlewood functions in $2m-1$ letters, in which the letter $1$ occurs once. Let
\begin{gather*}
\mu^-_x[f]=\sum_{\ell\in\Z}(-1)^\ell q^{L\ell^2/2}x^{L\ell}f(xq^\ell),\\
\Phi^+_a(x)=\mu^-_x[y^a+y^{-a}],\qquad \Psi^+_a(x)=\mu^-_x\Bigl[\frac{y^a+y^{-a}}{1+w^2+w(y+y^{-1})}\Bigr],
\end{gather*}
let $\delta^+_\varepsilon(x)=\det\bigl(x_i^{j-1+\varepsilon}+x_i^{-(j-1+\varepsilon)}\bigr)_{i,j=1}^m$ for $\varepsilon\in\{0,\frac12\}$, and
\[
P^B_t(\hat x)=\frac{(-wt;q)_\infty\prod_{i<m}(-wtx_i,-wt/x_i;q)_\infty}{(q;q)_\infty^m\prod_{i<m}(qx_i^{\pm2},qx_i^{\pm1};q)_\infty\prod_{1\le i<j<m}(qx_i^{\pm}x_j^{\pm};q)_\infty},
\]
where $(qx_i^{\pm2};q)_\infty=(qx_i^2,qx_i^{-2};q)_\infty$, and $(qx_i^\pm x_j^\pm;q)_\infty$ denotes the product over the four choices of signs. The frozen letter contributes the single factor $(-wt;q)_\infty$ and the cross factors $(qx_i^{\pm1};q)_\infty$.

\begin{proposition}[Even boundary formulas]\label{prop:evenbd}
\begin{enumerate}
\item For rational $c$,
\[
F_{2m,(0,\dots,0,c),0}(q)=(-q^{c+1/2};q)_\infty\,\HL^B_m(q^c;\one;q).
\]
\item For generic $\hat x$,
\begin{align*}
\HL^B_m(1;\hat x;q)&=P^B_1(\hat x)\,\frac{\det(\Phi^+_{j-1/2}(x_i))_{i,j=1}^m}{\delta^+_{1/2}(x)},\\
\HL^B_m(q;\hat x;q)&=(1+w)\,P^B_1(\hat x)\,\frac{\det(\Psi^+_{j-1/2}(x_i))_{i,j=1}^m}{\delta^+_{1/2}(x)},\\
\HL^B_m(w;\hat x;q)&=P^B_w(\hat x)\,\frac{\det(\Phi^+_{j-1}(x_i))_{i,j=1}^m}{\delta^+_{0}(x)}.
\end{align*}
\end{enumerate}
\end{proposition}

\begin{proof}[Proof sketch]
Part (i) is the case $n=2m-1$, $u=1$ of \cite[Theorem~3.7]{BW15}, after Euler's identity on the last coordinate as in \cref{lem:boundary} (equivalently \cite[Lemma~2.1]{CTW26}). For (ii), take level one, let $b\to\infty$ and put $c=-w/t$ in \cite[(5.6b)]{BW15}, which has $x=(\hat x,1)$, $K=n+\frac12$ and the factor $\Delta_B(-xq^r)/\Delta_B(-x)$. The one-variable factor becomes
\[
(-1)^{r}\,t^{r}\,q^{(m+1)r^2+(\frac12-m)r}\,x^{Lr}\,\frac{(-wx/t;q)_r}{(-wtx;q)_r}.
\]
Write $\Delta_B(y)=\prod_i(1-y_i)\prod_{i<j}(y_i-y_j)(y_iy_j-1)$ as $(-1)^m\prod_iy_i^{m-1/2}$ times the type-B determinant $\det\bigl(y_i^{m-j+1/2}-y_i^{-(m-j+1/2)}\bigr)$. At $y=-xq^r$ the monomial $\prod_i(-x_iq^{r_i})^{m-1/2}$ absorbs the factor $q^{(m-1/2)r}$, and for half-integral $a$ one has $(-y)^a-(-y)^{-a}=i^{2a}(y^a+y^{-a})$. Multilinearity in the rows turns the sum over $r\in\Z^m$ into the cosine determinant with entries $\Phi^+$; the constants cancel against the same computation at $r=0$, which produces $\delta^+_{1/2}$. This gives the case $t=1$. At $t=q$, the two telescoping factors combine with $P^B_q/P^B_1$ as in the proof of \cref{lem:boundary}; for the frozen letter the combination leaves the constant $1+w$. At $t=w$, the quotient $(1+x)/(1+xq^\ell)$ cancels a factor of the type-B denominator row by row, and the remaining determinant has integral indices.
\end{proof}

\begin{lemma}[Type-B tridiagonal inverse]\label{lem:evenB}
For $0\le k\le m$ and $\tau=\tau^{(0)}_k$ put
\[
J^B_k(\hat x)=P^B_1(\hat x)\,\frac{\det(\Phi^+_{\tau_j}(x_i))_{i,j=1}^m}{\delta^+_{1/2}(x)}.
\]
Then
\[
\sum_{j=0}^kA_{k,j}(q)\,\HL^B_m(q^j;\hat x;q)=q^{k^2/2}J^B_k(\hat x)\qquad(0\le k\le m).
\]
In particular $\HL^B_m(1;\hat x;q)-\HL^B_m(q;\hat x;q)=wJ^B_1(\hat x)$.
\end{lemma}

\begin{proof}
Put $d=m+1$ and index the functions by $a\in\{\frac12,\frac32,\dots,d-\frac12\}$. As in \cref{lem:bdexp}, $\Phi^+_a=(1+t^2)\Psi^{+[1]}_a+t(\Psi^{+[1]}_{a-1}+\Psi^{+[1]}_{a+1})$ with $t=w$, and similarly at every level $t=wq^{j-1}$. Two boundary effects change. At the bottom, the cosine symmetry gives $\Psi^+_{-1/2}=\Psi^+_{1/2}$, so the first diagonal entry of the tridiagonal matrix becomes $1+t+t^2$ and its determinant becomes $1+t$. At the top, the reindexing $\mu^-_x[f(y)]=-q^{L/2}\mu^-_x[y^Lf(qy)]$ applied to $f(y)=y^{-A}/(1+t/y)$ with $A=d+\frac12=(L+1)/2$ gives $\Psi^+_{d+1/2}+t\Psi^+_{d-1/2}=0$, so the last diagonal entry is $1$ as before. For $t=w$, the complementary-minor formula and Cauchy--Binet give
\[
\det\bigl(\Psi^+_{1/2},\dots,\Psi^+_{m-1/2}\bigr)=\frac{\det(\Phi^+_{1/2},\dots,\Phi^+_{m-1/2})-w\det(\Phi^+_{1/2},\dots,\Phi^+_{m-3/2},\Phi^+_{m+1/2})}{1+w},
\]
and \cref{prop:evenbd} gives the case $k=1$. For general $k$, the products $V_j$ of the modified matrices have determinant $(-w;q)_j$, which cancels in the complementary-minor formula, so \eqref{eq:bdexp} holds with the same coefficients $(V_j)_{d-h,d}$. For $j\le m$ no matrix path from the last coordinate reaches the modified first diagonal entry, so these coefficients agree with \cref{prop:coeffs}, and \cref{lem:inverse} applies verbatim.
\end{proof}

At $\hat x=\one$ the confluent limit is computed as in \cref{lem:confluent}, with the Vandermonde determinant in the variables $\log^2x_i$: since $\Phi^+_a(x)=\sum_{v\in a+L\Z}\epsilon_a(v)q^{(v^2-a^2)/(2L)}(x^v+x^{-v})$,
\begin{equation}\label{eq:evenconf}
J^B_k(\one)=\frac{(-w;q)_\infty^{2m-1}}{(q;q)_\infty^\nu}\sum_{v\in\tau+L\Z^m}\epsilon_\tau(v)\frac{\chi_0(v)}{\chi_0(\tau^{(0)}_*)}q^{\frac{\|v\|^2-\|\tau\|^2}{2L}},\qquad\tau=\tau^{(0)}_k.
\end{equation}

\begin{theorem}[Even companions in sector (i)]\label{thm:evencomp}
Let $m\ge2$ and $1\le k\le\min(6,m-1)$. Then the identity of \cref{conj:evenwron} holds for $B^{\mathrm{(i)}}_k$, that is, $F_{2m,B^{\mathrm{(i)}}_k,0}=R^{(0)}_{\tau^{(0)}_k}$. Consequently $\bigl(B^{\mathrm{(i)}}_k,C\bigr)\in\M_{2m}$ with the scalar $C$ of \cref{conj:even}(i).
\end{theorem}

\begin{proof}
The vector $B^{\mathrm{(i)}}_k$ is zero except for $k$ pairs $(-1,1)$ in positions $(2m-2i,2m-2i+1)$, $1\le i\le k$. For $k=1$, the contiguity computation leading to \eqref{eq:contig2} applies verbatim, since $T_{2m}(e_{2m-1}+e_{2m})=e_{2m-1}-e_{2m-2}$; with $c=0$ and \cref{prop:evenbd}(i) it gives $F_{2m,B^{\mathrm{(i)}}_1,0}=q^{-1/2}(-w;q)_\infty(\HL^B_m(1;\one;q)-\HL^B_m(q;\one;q))$. For $2\le k\le6$, the reductions of \cref{lem:k2red,lem:certred} at $v=0$, that is $u=1$, involve only the last $2k+1$ coordinates, and their proofs use only the local form of the tadpole quadratic form near the last coordinate. They therefore hold in rank $2m$ whenever $2m\ge2k+1$, and give
\[
q^{k^2/2}F_{2m,B^{\mathrm{(i)}}_k,0}=\sum_{j=0}^kA_{k,j}(q)\,(-w;q)_j\,F_{2m,je_{2m},0}=(-w;q)_\infty\sum_{j=0}^kA_{k,j}(q)\,\HL^B_m(q^j;\one;q).
\]
(For $k=5,6$ this uses the certificates at $u=1$.) \Cref{lem:evenB} and \eqref{eq:evenconf} identify the right-hand side with $R^{(0)}_{\tau^{(0)}_k}$, and \cref{prop:evenwron} gives modularity.
\end{proof}

\begin{corollary}\label{cor:r6more}
The second and third pairs of \cref{thm:r6} are modular. Thus four of the nine pairs in rank six are proved modular, and for every $m\ge7$ the set $\M_{2m}$ contains at least the eight pairs of \cref{conj:even}(i) with $k\le6$ and (ii) with $j=0$.
\end{corollary}

The remaining even sectors, namely (i) with $k=m$, (ii) with $j\ge1$, (iii) and (iv), contain the shift $v=\frac12$ of the boundary family $B=(2v,-2v,\dots,2v,c)$. For $v\ne0$, \cite[Theorem~3.7]{BW15} gives Hall--Littlewood functions in the letters $(u,u^{-1},\dots,u^{-1},u)$ with $u=q^v$, in which the unpaired letter is $u$ rather than $1$. This is not of the form covered by \cite[(5.6b)]{BW15}, and a determinant formula for it is the missing ingredient for these sectors. In plethystic terms the obstruction is the unpaired letter: for $u=w$ one has $p_k[X/(1-q)]=\bigl((m-1)(w^k+w^{-k})+w^k\bigr)/(1-q^k)$, which is not of the form produced by the alphabets $(x^\pm)$ or $(\hat x^\pm,1)$.

\begin{remark}\label{rem:evenlead}
The constant terms of the right-hand sides in \cref{conj:evenwron} are
\[
\frac{\chi_0(\tau^{(0)}_k)}{\chi_0(\tau^{(0)}_*)}=\binom{2m-1}{k}-\binom{2m-1}{k-2},\qquad c_j\,\frac{\chi_0(\tau^{(1/2)}_j)}{\chi_0(\tau^{(1/2)}_*)}=\binom{2m}{j},
\]
by cancellation in the products defining $\chi_0$. These are the even analogues of \cref{prop:leading}. The factor $c_m=2$ is exactly what is needed for $\binom{2m}{m}$. The Nahm sums have the same constant terms in every case computed (ranks two to fourteen), and a path-counting proof as in \cref{prop:leading} should be possible.
\end{remark}


\begin{remark}\label{rem:evenwron-evidence}
The identities of \cref{conj:evenwron} not covered by \cref{thm:evencomp} or \cite{CTW26} agree with the Nahm sums as $q$-series, in exact integer arithmetic, through $q^{40}$ in rank six, through $q^{15}$ in ranks eight, ten and twelve, and through $q^{10}$ in rank fourteen. The relation for $B^{\mathrm{(iv)}}$ agrees through $q^{20}$ in ranks four to eight. In rank four the corresponding sums are known to be modular by \cite{SW26,SW26b}. This evidence is not used in any proof.
\end{remark}

\begin{remark}\label{rem:highrank}
The conjectures have been tested against the cusp conditions beyond the classified ranks. For each of the $11$ and $13$ pairs predicted by \cref{conj:even} in ranks eight and ten, and for each of the $18$ and $21$ pairs predicted by \cref{conj:odd} in ranks eleven and thirteen, the procedure of \cref{sec:cusp} gives $\kappa_2=\kappa_3=0$ and $c_1-\gamma_0=C$ to within $10^{-10}$ in floating-point arithmetic. Random half-integral shifts give values between $10^{-2}$ and $10^{2}$. In rank eleven, replacing the ordered chain $x_1\le\cdots\le x_5$ by unordered patterns such as $(1,0,0,0,0)$ or $(1,0,1,0,1)$ gives $|\kappa_2|>2\cdot10^{-2}$; thus the ordering in \cref{conj:odd} is not an artefact of the ranks used to formulate it. These computations are evidence only.
\end{remark}

\begin{remark}\label{rem:numerics}
As a check on normalizations only, identity \eqref{eq:wron} has been compared coefficientwise, in exact integer arithmetic, through $q^{15}$ for $(m,k)=(8,7)$ and $(9,8)$ in both sectors and for $(m,k)=(9,7)$ and $(10,9)$ in the sector $z=0$. These cases lie outside the proved range. Such comparisons are not used anywhere as proofs.
\end{remark}

\subsection{The missing boundary reduction}
Set $\HL_{m,j}=\HL_m(q^j;\one;q)$. The reduction in the sector $z=0$ is
\begin{equation}\label{eq:missing}
\frac{S^{(0)}_{m,k}}{(-w;q)_\infty}=q^{-k^2/2}\sum_{j=0}^kA_{k,j}(q)\,\HL_{m,j}.
\end{equation}
\Cref{thm:basic,thm:k2,thm:k3to6}, together with \cref{lem:inverse}, prove this identity for $0\le k\le\min(6,m)$ and for $k=m$. The remaining interior indices are conjectural. The local formulation below would prove the more general reduction \eqref{eq:certred} by induction.

\subsection{A local induction problem for the remaining companions}
The boundary relation at index $m+1$ supplies the initial value of a more demanding contiguity identity, which we now state explicitly; the basis theorem alone does not prove it.

Retain $u=q^v$, $w=q^{1/2}$ and the polynomials $R_s(e)$ defined above, now for every $s\ge0$. For $k\ge1$, put $r_k=2k-1$, and let $b^{(k)}_v\in\Q^{r_k}$ consist of $k-1$ pairs $(-2v,2v)$ followed by $-v$. For integers $N,s\ge0$, define
\begin{equation}\label{eq:moment}
\mathfrak M^{(s)}_{k,N}(v;q)=\sum_{j=0}^sA_{s,j}(q)\,(-w/u;q)_j\,F_{r_k,\,b^{(k)}_v-Ne_1+je_{r_k},\,0}(q).
\end{equation}
For $k\ge2$, \cref{cor:bdrel} proves $\mathfrak M^{(k)}_{k,0}=0$, and Euler's identity proves the same for $k=1$. The proposed extension is
\begin{equation}\label{eq:extension}
\mathfrak M^{(k)}_{k,N}(v;q)=\frac{q^{k-1/2}}u\,(q^{-N}-1)\,\mathfrak M^{(k-1)}_{k,N-1}(v;q)\qquad(N\ge1).
\end{equation}
Both moments here have rank $2k-1$. We do not establish this identity for arbitrary $k$.

The following proposition gives its precise role. For $n=(n_1,\dots,n_{2k})\in\Z_{\ge0}^{2k}$ and indeterminates $a=(a_1,\dots,a_{2k})$, set
\[
W_k(n;u)=\frac{q^{\sum_{i=1}^{2k-1}n_i^2+n_{2k}^2/2-\sum_{i=1}^{2k-1}n_in_{i+1}}\,u^{2\sum_{i=1}^{2k-1}(-1)^{i+1}n_i-n_{2k}}}{\prod_{i=1}^{2k}(q;q)_{n_i}},\qquad T_k(a)=\frac{a_1a_3\cdots a_{2k-1}}{a_2a_4\cdots a_{2k-2}},
\]
where an empty product is $1$. For a Laurent polynomial $P$ in $2k$ variables, write
\[
\mathcal Z_k(X;P)=\sum_{n\ge0}W_k(n;u)\,X^{n_1}P(q^{n_1},\dots,q^{n_{2k}}).
\]
Positive definiteness of $T_{2k}$ makes this an entire function of $X$. The local boundary identity is
\begin{equation}\label{eq:local}
\mathcal Z_k(X;R_k(a_{2k}))=q^{k^2/2}u^kX\,\mathcal Z_k(X;T_k).
\end{equation}

\begin{proposition}[Induction criterion]\label{prop:induction}
Identity \eqref{eq:local} holds for $k=1$ and every rational $v$. For $k\ge2$, assume that it holds at $k-1$ for every $X$. Then \eqref{eq:local} at $k$ is equivalent to \eqref{eq:extension} for every $N\ge1$. If these identities hold through index $k$, they imply \eqref{eq:certred} at that index for every $m\ge k$.
\end{proposition}

\begin{proof}
The coefficient of $X^N$ on the left of \eqref{eq:local} is
\[
\frac{q^{N^2}u^{2N}}{(q;q)_N}\,\mathfrak M^{(k)}_{k,N}.
\]
For $N\ge1$, fix $n_1=N-1$ on the right. The remaining insertion is the full companion of index $k-1$ on the coordinates $n_2,\dots,n_{2k}$, with an additional first-coordinate shift $-(N-1)e_1$. Apply the local identity at $k-1$, first fixing $n_2$ and then summing over it. The resulting coefficient is
\[
\frac{q^{N^2-N+k-1/2}u^{2N-1}}{(q;q)_{N-1}}\,\mathfrak M^{(k-1)}_{k,N-1}.
\]
Equality of these coefficients is precisely \eqref{eq:extension}. The constant coefficients vanish on both sides, since $\mathfrak M^{(k)}_{k,0}=0$. Since both sides are entire in $X$, this proves the equivalence.

For the base case, put $c=w/u$. Euler's identity gives
\[
\mathfrak M^{(0)}_{1,N}=(-cq^{-N};q)_\infty,\qquad \mathfrak M^{(1)}_{1,N}=c(q^{-N}-1)(-cq^{1-N};q)_\infty,
\]
and substitution into the coefficient calculation verifies \eqref{eq:local} for $k=1$, including its constant term.

Finally, fix the coordinate $n_{r-2k}$ of a Nahm sum of rank $r=2m+1$ and set $X=q^{-n_{r-2k}}$. The factor $XT_k$ is exactly the insertion for its last $k$ companion pairs. Summing \eqref{eq:local} over the fixed coordinate and all preceding coordinates, which is justified by absolute convergence, yields
\[
q^{kv+k^2/2}G_{m,k}(v;q)=\sum_{j=0}^kA_{k,j}(q)\,(-w/u;q)_j\,F_{r,B_v+je_r,0}(q).
\]
\Cref{lem:boundary} then gives \eqref{eq:certred}.
\end{proof}

The telescoping certificates of \cref{sec:cert} prove the local identity through $k=4$ for arbitrary rational $v$, and through $k=6$ for $v=0,\frac12$; consequently they also prove \eqref{eq:extension} in those ranges. For larger indices, the value at $N=0$ alone does not determine the values at $N\ge1$, and establishing that extension is the remaining step in the proposed induction. \Cref{prop:recurrence} is an identity among the polynomials $R_k$ and does not supply this extension.
